\scanpage{23}

\margin{\tikz[baseline=(r.base)]{\node[inner sep=1pt] (r) {$\mathfrak h\subset M\subset M_a$};\draw (r.north west)--(r.south west)--(r.south east);}}

% Three diagonal strokes run through the draft. The dots at the right are ink.
{\scriptsize
\diagram{\begin{tikzcd}[ampersand replacement=\&,row sep=2pt,column sep=2pt,cells={anchor=west}]
{}\arrow[from=1-1,to=6-1,draw=none,"{\left\{\vphantom{\rule{0pt}{4.4em}}\right.}"{left,inner sep=1pt}]
\arrow[from=9-1,to=12-1,draw=none,"{\left\{\vphantom{\rule{0pt}{3.7em}}\right.}"{left,inner sep=1pt}]
\arrow[dash,from=1-2,to=12-2,start anchor=north west,end anchor=south]
\arrow[dash,from=1-2,to=10-2,start anchor=north,end anchor=south east] \&\deleted{(\mathrm{i})\ ,\ \mathfrak h\text{ contient une sous-algèbre de Cartan}}\\
{} \&\deleted{(\mathrm{ii})\quad\mathfrak h=M,\quad\mathfrak h\text{ contient \unsure{élémt|élément|des éléments} régulier de }\mathfrak g}\\
{}\arrow[ddd,Rightarrow] \&\deleted{(\mathrm{iii})\quad\exists\,a\in\mathfrak h\quad\text{t.q. }\operatorname{ad}_{\mathfrak g/\mathfrak h}(a)\text{ injectif}\quad
(\text{i.e. tel que }\deleted{\unsure{M_a|M}}\mathfrak h=M_a)}\\
{} \&\deleted{(\mathrm{iv})\quad\psi:X\unsure{\mathbin{-}|\longrightarrow}W(\mathfrak g)\quad\text{dominant, et }\mathfrak h=M\qquad\bullet}\\
{} \&\deleted{(\mathrm{v})\quad\psi:X\mathbin{-}W(\mathfrak g)\quad\text{gén. lisse}}\\
{} \&\deleted{(\mathrm{vi})\quad\varphi:X\times W(\mathfrak h)\longrightarrow W(\mathfrak g)\quad\text{gén\textsuperscript{t} lisse}}\\
{} \&\deleted{(\mathrm{vii})\quad\psi:X\longrightarrow W(\mathfrak g)\quad\text{dominant}\qquad\bullet}\\
{}\arrow[u,Rightarrow] \&\deleted{(\mathrm{viii})\quad\psi:X\mathbin{-}W(\mathfrak g)\quad\text{dominant et quasi-fini}\qquad\bullet}\\
{}\arrow[u,Rightarrow] \&\deleted{(\mathrm{ix})\quad\psi\quad\text{gén. \unsure{qu.|quasi} fini et }\mathfrak h=M}\\
{}\arrow[uuuuuuu,Rightarrow,bend left=25] \&\deleted{\deleted{(\unsure{\mathrm{viii}|\mathrm{xiii}})}(\mathrm{x})\qquad\deleted{\psi\text{ dominant}}\quad\psi:X\longrightarrow W(\mathfrak g)\quad\text{dominant, \uline{et}}\quad\bullet}\\
{} \&\deleted{\left[\begin{aligned}&\mathfrak h\text{ est l'alg. de Lie d'un sous-groupe lisse }H\\&\text{de }\mathfrak g\end{aligned}\right.}\\
{} \&\deleted{\deleted{(\unsure{\mathrm{x}|\mathrm{ix}|\mathrm{xii}})}(\mathrm{xi})\quad\psi:X\mathbin{-}W(\mathfrak g)\quad\text{dominant, et }N\text{ lisse.}\qquad\bullet}
\end{tikzcd}}}

\deleted{$\mathfrak h$ \deleted{\unsure{plus}} minimale parmi les $k$. algèbres de Lie qui sont}

\deleted{a) soit $={}$ normalisateur, et \unsure{satisfont|satisfait|satisfasse} \unsure{au|aux} th. de densité}

\deleted{b) soit $={}$ normalisateur, et satisfont aux th. des \unsure{dérivations|désintégrations}}

\deleted{c) soit définies par $H$ lisse, et \rule{5em}{0.3pt} densité}

\deleted{d) soit telles que $\exists\,a\in\mathfrak h\quad\operatorname{ad}(a)_{\mathfrak g/\mathfrak h}$ injectif}

\deleted{e) soit telles que}

% archive sheet number: 22
